BKS chiude il percorso passando dal modello parametrico alla tabella empirica dei comportamenti. Con $L$ strumenti binari esistono al massimo $2^L$ pattern possibili. Per ciascun pattern $E$ e per ciascuna famiglia $f$, BKS conta quante volte quel pattern e' stato associato a una classe nel set di calibrazione:
\begin{equation}
N(E,\Vuln),\qquad N(E,\Safe).
\end{equation}
La stima non smussata sarebbe
\begin{equation}
P(\Vuln\mid E)=
\frac{N(E,\Vuln)}{N(E,\Vuln)+N(E,\Safe)}.
\end{equation}
La repository usa pero' un m-estimate, che evita stime estreme quando un pattern e' raro:
\begin{equation}
s_{\mathrm{BKS}}(x)=
\frac{N(E,\Vuln)+m\pi_f}{N(E,\Vuln)+N(E,\Safe)+m},
\end{equation}
dove $\pi_f$ e' il prior vulnerabile della famiglia e $m$ controlla quanto lo score viene riportato verso il prior.

Nel nostro esempio, il pattern $(1,0,0)$ compare $50$ volte: $42$ vulnerabili e $8$ safe. Con $\pi_f=0.40$ e $m=2$,
\begin{equation}
s_{\mathrm{BKS}}(x^\star)=
\frac{42+2\cdot 0.40}{50+2}
\approx 0.82.
\end{equation}
La decisione e' vulnerabile. BKS cattura direttamente il fatto che ``CodeQL segnala mentre Semgrep e Joern tacciono'' puo' essere un pattern storicamente forte, anche se majority voting, voto pesato e Naive Bayes lo classificano safe. Questa e' la sua forza principale: non impone additivita' ne' indipendenza condizionata. Il prezzo e' la sparsita': al crescere degli strumenti, molti dei $2^L$ pattern saranno rari o assenti, quindi servono calibrazione sufficiente e smoothing per evitare overfitting.

\begin{lstlisting}[caption={Implementazione BKS in analysis/fusion/bks.py}]
from analysis.fusion.bks import bks_predictions

preds = bks_predictions(
    calibration_df,
    validation_df,
    tools,
    families,
    labels,
    fire_index,
    tau=0.5,
    m=2.0,
)
\end{lstlisting}
